% !TEX root = ../main.tex
\section{Visão Geral do Projeto}

\subsection{Descrição do Projeto}

O Educado é uma plataforma educacional gamificada e mobile-first, criada para oferecer cursos curtos e acessíveis a catadores de materiais recicláveis.
O projeto é desenvolvido por uma colaboração internacional entre universidades e organizações de quatro países, com participação de catadores no processo de criação da solução.

A plataforma atende três grupos principais: os catadores, que acessam os conteúdos educacionais; as equipes universitárias, responsáveis pela criação e curadoria dos cursos; e os administradores, que gerenciam usuários, instituições e informações da plataforma.
Entre suas funcionalidades estão a organização de cursos em seções e atividades, o acompanhamento do progresso dos estudantes, a gamificação, a emissão de certificados e a apresentação de indicadores de uso.

O projeto está distribuído em quatro repositórios públicos \cite{educado-repos}: educado-web, educado-api, educado-app e educado-docs.
Por recomendação da disciplina, este trabalho tem como objeto principal o educado-web, responsável pela interface web utilizada por criadores de conteúdo e administradores.
A caracterização da estratégia de testes, porém, cobre também o educado-api, o educado-app e o educado-docs, porque a prática de teste do projeto se distribui entre eles e analisar apenas a interface web daria um retrato incompleto.

\subsection{Características Técnicas}

O Educado é organizado em quatro repositórios, cada um responsável por uma parte da plataforma.
Essa separação permite que o backend, as interfaces de usuário e a documentação sejam desenvolvidos e implantados de forma independente.

\begin{table}[H]
\centering
\caption{Repositórios que compõem o projeto Educado.}
\label{tab:repositorios}
\begin{tabularx}{\textwidth}{l X X}
\toprule
\textbf{Repositório} & \textbf{Responsabilidade} & \textbf{Tecnologias principais} \\
\midrule
educado-web & Interface web para criadores de conteúdo e administradores & TypeScript, Vite, CSS e Vitest \\
educado-api & API REST, regras de negócio e persistência dos dados & Node.js, Express, PostgreSQL e Jest \\
educado-app & Aplicativo móvel utilizado pelos estudantes & React Native, Expo e Jest \\
educado-docs & Documentação do produto e de sua arquitetura & MkDocs e Python \\
\bottomrule
\end{tabularx}
\end{table}

O educado-web, foco deste relatório, é uma aplicação de página única (Single Page Application, SPA) desenvolvida em TypeScript, sem a utilização de frameworks de interface como React, Vue ou Angular.
A renderização das páginas é realizada diretamente por funções que manipulam o HTML e registram os eventos necessários.
O projeto também possui implementações próprias para o roteamento da aplicação e para a internacionalização das interfaces em português e inglês.

A aplicação utiliza o Vite como servidor de desenvolvimento e ferramenta de compilação.
Os estilos são escritos em CSS puro e organizados de acordo com as funcionalidades do sistema.
Entre as dependências de execução estão o Axios e o Flatpickr, embora o cliente HTTP utilizado atualmente seja baseado na API nativa fetch.
O TypeScript está configurado em modo estrito, incluindo verificações de parâmetros, variáveis não utilizadas e fluxos incompletos em estruturas condicionais.

O código-fonte segue uma organização modular por funcionalidades.
Os módulos de autenticação, cursos, mídia, administração, dashboard, páginas públicas e interface do estudante encontram-se no diretório \texttt{src/features}.
Os recursos compartilhados, como comunicação HTTP, sessão de autenticação, internacionalização e componentes de interface, estão concentrados em \texttt{src/shared}.
O roteamento e os estilos globais ficam em \texttt{src/app}.

O frontend não possui banco de dados próprio.
A persistência e as regras de negócio são fornecidas pelo educado-api, acessado por meio de requisições HTTP.
A URL do serviço é configurada pela variável de ambiente \texttt{VITE\_API\_URL}, enquanto a autenticação utiliza um token JWT fornecido pela API e armazenado no \texttt{localStorage} do navegador.
Essa dependência é relevante para o teste: verificar a interface em conjunto com o backend exigiria disponibilizar a API, o que a estratégia atual do projeto não faz.

Para execução local, o projeto requer Node.js 20 ou superior e o gerenciador de pacotes npm.
A instalação reprodutível das dependências é feita por meio do arquivo \texttt{package-lock.json}.
Em produção, a aplicação pode ser compilada em uma imagem Docker de múltiplos estágios e disponibilizada por um servidor Nginx, configurado para atender às rotas da SPA e armazenar recursos estáticos em cache.

Na versão analisada, o diretório \texttt{src} possui aproximadamente 81 arquivos TypeScript e CSS, totalizando cerca de 34 mil linhas.
A maior parte do código está escrita em TypeScript, com aproximadamente 21 mil linhas, enquanto os arquivos de estilo somam aproximadamente 14 mil linhas.
Esses valores são aproximados e incluem comentários, linhas em branco e dados de internacionalização.

\subsection{Processo de Desenvolvimento}

O Educado é desenvolvido de forma colaborativa por equipes vinculadas a universidades e organizações de diferentes países.
Segundo a proposta do projeto, o trabalho utiliza práticas ágeis, com ciclos quinzenais, revisões internacionais e comunicação contínua entre as equipes.
O código-fonte é mantido publicamente no GitHub sob a licença Apache 2.0.

No \texttt{educado-web}, a documentação de contribuição estabelece a branch \texttt{main} para versões estáveis e a branch \texttt{dev} para integração das alterações.
Novas contribuições devem partir da \texttt{dev} e utilizar branches identificadas por prefixos como \texttt{feat/}, \texttt{fix/}, \texttt{docs/} e \texttt{chore/}.
Os commits seguem a convenção \textit{Conventional Commits}, facilitando a identificação do tipo e do propósito de cada mudança.

As contribuições são submetidas por meio de \textit{pull requests}.
O modelo disponibilizado pelo projeto solicita a descrição da alteração, a associação com issues, os procedimentos utilizados para testá-la e, no caso de mudanças visuais, capturas de tela ou gravações.
O checklist também inclui verificações de compilação, internacionalização, navegação, comunicação com a API e tratamento seguro dos dados inseridos na interface.

O repositório utiliza GitHub Actions para executar automaticamente, a cada \textit{pull request}, a instalação das dependências, a verificação estática do TypeScript, os testes automatizados e a compilação da aplicação.
O projeto também fornece modelos específicos para o registro de defeitos e a solicitação de funcionalidades, além dos documentos \texttt{README.md}, \texttt{CONTRIBUTING.md}, \texttt{CODE\_OF\_CONDUCT.md} e \texttt{SECURITY.md}.
Esses artefatos orientam a configuração do ambiente, a contribuição de código e a comunicação de problemas técnicos ou de segurança.
